\section{Infinite-dimensional calculus and manifolds}\label{App:Bast}

In this appendix, we recall some basic definitions of infinite-dimensional calculus.
As stated in the beginning, we will mostly work on Banach manifolds, but the manifolds of mappings cannot be modelled on Banach spaces. Hence we need to base our article on a generalised calculus, the so called Bastiani calculus \cite{Bas64}, see, e.g., \cite{GaS22, Sch23} for introductory expositions.

\begin{Definition}%\label{defn-Ck}
Consider locally convex spaces $E$, $F$
and a map $f\colon U\rightarrow F$ on an open subset $U\subseteq E$.
Write
\[
D_yf(x):= {\rm d}f(x;y):= \frac{{\rm d}}{{\rm d}t}\biggr|_{t=0}f(x+ty)
\]
for the directional derivative of~$f$ at $x\in U$
in the direction $y\in E$, if it exists.
Let $k\in \N_0\cup\{\infty\}$.
If $f$ is continuous, and the iterated directional derivatives
\[
{\rm d}^jf(x;y_1,\dots, y_j):= (D_{y_j}\cdots D_{y_1}f)(x)
\]
exist for all $j\in\N_0$ such that
$j\leq k$, $x\in U$ and $y_1,\dots, y_j\in E$,
and the maps ${\rm d}^jf\colon U\times E^j\to F$
are continuous, then $f$ is called~$C^k$.
A map of class $C^\infty$ is also called~\emph{smooth}.
\end{Definition}

The chain rule holds in the usual form for Bastiani differentiable mappings (see, e.g., \cite[Proposition~1.23]{Sch23}), whence it makes sense to define manifolds as in the finite-dimensional setting via charts. Hence a continuous mapping $f \colon M \rightarrow N$ between (possibly infinite-dimensional, smooth) manifolds is smooth if, locally in suitable pairs of charts the map is Bastiani smooth.

%\begin{setup}%\label{conventions-mfd}
All manifolds and Lie groups considered
in the article are modeled on locally
convex spaces which may be infinite-dimensional,
unless the contrary is stated. If we say that a manifold or a~vector bundle is a Banach manifold or a Banach vector bundle we mean that all modelling spaces (or fibres in the bundle case) are Banach spaces. See~\cite{La01} for a discussion of differential geometry on Banach manifolds. All manifolds are assumed to be smooth and Hausdorff.
We~write~$\pi_M \colon TM \rightarrow M$ for the tangent bundle with bundle projection $\pi_M$.
%\end{setup}

In the Bastiani setting, the notion of submersion has to be defined as in \cite[Definition 4.4.8]{Ham82}, see also \cite{Sch23}.

\begin{Definition}\label{defn:submersion}
Let $f \colon M \rightarrow N$ be a smooth map between manifolds, we say that $f$ is a~\emph{submersion}, if for every $x \in M$ there exist manifold charts of $M$ around $x$ and of $N$ around $f(x)$ which conjugate the map $f$ to a projection onto a complemented subspace of the model space of~$M$. Pairs of charts with these properties will be called \emph{submersion charts} for the submersion~$f$.\looseness=1
\end{Definition}
We refer to \cite{glo16} for a discussion of the notion of submersion in the infinite-dimensional setting and its properties.
In particular, with this definition of a submersion fibre-products of manifolds can be defined as in the finite-dimensional setting. In particular, the usual construction for a~submersion $p\colon M \rightarrow N$ and a smooth map $g\colon L \rightarrow M$ yields (see \cite[Lemma 1.60]{Sch23}) that the fibre-product
\begin{displaymath}
 M {}_{p}\hspace{-.25em}\times_{g} L := \{(x,y) \in M\times M \mid p(x)=g(y)\}
\end{displaymath}
is a submanifold of $M\times M$. If $g=p$, we also write shorter $M\times_p M$ for this fibre-product.

\begin{Lemma}\label{lem:tangent_fibreprod}
Let $p_i \colon M_i \rightarrow N$, $i=1,2$, be surjective submersions, then the identification $T(M_1\times M_2) \cong TM_1 \times TM_2$ restricts to a diffeomorphism
\begin{displaymath}
 T((M_1) {}_{p_1}\hspace{-.5em}\times_{p_2} (M_2) ) \cong (TM_1) {}_{Tp_1}\hspace{-.5em}\times_{Tp_2} (TM_2).\end{displaymath}
\end{Lemma}

\begin{proof}
This question is clearly local, so we pick submersion charts on $W_i \subseteq M_i$, $i=1,2$, for~$p_i$ such that $W_i \cong U \times V_i$, with $U\subseteq E$, $V_i \subseteq F_i$ open in locally convex spaces and $p_i$ becomes the projection onto the first component with respect to this splitting. Suppressing the submersion charts, the embedding $M_1\times_N M_2 \subseteq M_1 \times M_2$ becomes
\begin{displaymath}
 \iota \colon\ U \times V_1 \times V_2 \rightarrow (U\times V_1) \times (U \times V_2), (u,v_1,v_2) \mapsto ((u,v_1),(u,v_2)).
\end{displaymath}
Passing to the tangent, we obtain the embedding $T\iota$
\begin{align*}
 U \times V_1 \times V_2 \times (E \times F_1) \times (E\times F_2) &\rightarrow ((U\times V_1)\times E\times F_1)\times ((U\times V_2)\times E\times F_2),\\
 (u,v_1,v_2,x,y_1,y_2)&\mapsto ((u,v_1,x,y_1),(u,v_2,x,y_2) ).
\end{align*}
Moreover, as $Tp_i|_{W_i} \colon (U \times V_i) \times E \times F_i \rightarrow U\times E, ((u,v_i),x,y_i) \mapsto (u,x)$, we see that as sets (hence also as manifolds) $T\iota (T((W_1){}_{p_1}\hspace{-.4em}\times_{p_2} \hspace{-.2em} (W_2)) ) ) =(TW_1){}_{Tp_1}\hspace{-.4em}\times_{Tp_2}\hspace{-.2em}(TW_2)$ as claimed.
\end{proof}
